\documentclass[11pt]{article}
\usepackage[letterpaper,margin=1in]{geometry}
\usepackage[T1]{fontenc}
\usepackage{lmodern,microtype}
\usepackage{amsmath,amssymb}
\usepackage{xcolor}
\usepackage[colorlinks=true,linkcolor=blue!40!black,citecolor=blue!40!black,urlcolor=blue!40!black]{hyperref}
\usepackage{fancyhdr}
\pagestyle{fancy}
\fancyhf{}
\fancyhead[L]{\small Bach flow and smooth Poincar\'e}
\fancyhead[R]{\small Working draft}
\fancyfoot[C]{\thepage}
\setlength{\headheight}{14pt}
\setlength{\emergencystretch}{2em}
\numberwithin{equation}{section}
\DeclareMathOperator{\Ric}{Ric}
\newcommand{\CP}{\mathbb{CP}}

\title{Bach flow and the smooth four-dimensional\\Poincar\'e conjecture}
\author{Edward B. Baker III}
\date{Working draft --- 23 November 2026}
% Prepared for Edward B. Baker III with OpenAI Codex assistance,
% 23 November 2026 (America/New_York). The selected model and effort were unavailable.
% The common opening is copied from the companion chiral metric evolution draft.
% The exact source identity and evidence scope are recorded in ../DRAFT_HISTORY.md.
% This draft contains the introduction and its cited background.

\begin{document}
\maketitle
\thispagestyle{plain}

\section{Introduction}\label{sec:introduction}

Isothermal coordinates provide a simple model for using geometry to
recognize the smooth structure of a sphere. Every smooth Riemannian metric
on an oriented real surface can be written locally as
\[
 g=e^{2u}(dx^2+dy^2)=e^{2u}|dz|^2,\qquad z=x+iy.
\]
Orientation-preserving transitions between these charts are holomorphic,
so the metric determines a Riemann surface. If the underlying smooth
surface is closed and homotopy equivalent to \(S^2\), it is orientable
and simply connected. Uniformization then identifies this compact Riemann
surface biholomorphically with \(\CP^1\); the resulting map is a
diffeomorphism to the standard smooth two-sphere~\cite{Donaldson2011}.
This gives a geometric proof of the two-dimensional smooth Poincar\'e
statement. Isothermal coordinates supply the local conformal description,
while uniformization supplies the global identification; the existence of
the local charts alone would not establish the conclusion.

In three dimensions, the analogous statement is that every closed smooth
three-manifold homotopy equivalent to \(S^3\) is diffeomorphic to the
standard three-sphere. Hamilton's approach replaces conformal coordinates
by a geometric heat equation, the Ricci flow,
\[
 \partial_t g=-2\Ric(g).
\]
After fixing the diffeomorphism freedom, this is a nonlinear parabolic
evolution of the metric. Hamilton proved that in dimension three the
volume-normalized flow starting from strictly positive Ricci curvature
converges to a metric of constant positive sectional
curvature~\cite{Hamilton1982}. General initial metrics can develop
singularities, so the broader program required both their analysis and
a way to continue the evolution. Perelman's entropy and noncollapsing
estimates supplied decisive control of singularity models~\cite{Perelman2002}.
His construction of Ricci flow with surgery removes appropriate neck
regions and replaces them by controlled caps~\cite{Perelman2003Surgery}.
For a homotopy three-sphere, finite-time extinction together with the
topology of these surgeries yields the smooth identification with
\(S^3\)~\cite{Perelman2003Extinction}. The lesson for higher dimensions is
that a local smoothing equation must be accompanied by global control of
degeneration and its topological consequences.

In higher dimensions, Smale's handle-cancellation methods, underlying
$h$-cobordism theory, show that every closed smooth homotopy $n$-sphere
with $n\geq5$ is homeomorphic to $S^n$~\cite{Smale1961}.
Milnor's exotic seven-spheres show that the same topological sphere can
admit distinct smooth structures~\cite{Milnor1956}.
Dimension four lies outside the range of these smooth handle-cancellation
arguments. Freedman proved that every closed topological homotopy
four-sphere is homeomorphic to \(S^4\)~\cite{Freedman1982}. The smooth
four-dimensional Poincar\'e conjecture asks whether every closed smooth
homotopy four-sphere is diffeomorphic to the standard \(S^4\). This remains
open. Major approaches include handle calculus and surgery descriptions
that seek to simplify homotopy
spheres, and knot-theoretic obstructions that could distinguish a candidate
exotic sphere~\cite{FreedmanGompfMorrisonWalker2010}.
Geometric approaches seek a distinguished metric or connection
whose existence would force the standard smooth structure. The present
work pursues this direction, with anti-self-dual (ASD) conformal geometry
as a possible target. Fine~\cite{Fine2011} introduced a gradient flow for
definite \(SO(3)\) connections aimed at producing ASD Einstein metrics of
nonzero scalar curvature. This offers a concrete
four-dimensional evolution framework; its application to the smooth
Poincar\'e problem would still require suitable initial connections on
arbitrary homotopy spheres and global control of the flow.

The special geometry of four dimensions explains the appeal of these
variables. Unlike a surface metric, a general four-dimensional metric has
Weyl curvature that cannot be removed by conformal rescaling or a change
of coordinates. At the same time, on an oriented Riemannian four-manifold
the Hodge star satisfies \(*^2=1\) on real two-forms, producing two real
rank-three bundles,
\[
 \Lambda^2T^*M=\Lambda_g^+\oplus\Lambda_g^-.
\]
Atiyah, Hitchin, and Singer~\cite{AHS1978} used this chiral splitting to
connect conformal geometry and gauge fields. In our orientation convention,
their twistor integrability criterion says that \(W^+=0\) precisely when
the natural almost complex
structure on the sphere bundle \(S(\Lambda_g^+)\) is integrable.
The resulting complex threefold encodes a real four-dimensional conformal
geometry through a family of projective lines. Penrose's nonlinear-graviton
construction supplies the corresponding
holomorphic picture and, with additional data, its Ricci-flat
specialization~\cite{Penrose1976}. These are four-dimensional analogues of
the relation between conformal geometry and complex analysis on surfaces,
with a curvature integrability condition in place of an automatic local
complex structure.

Gauge theory provides a complementary reason to work with connections.
Donaldson invariants extract information about smooth four-manifolds from
moduli spaces of anti-self-dual gauge connections~\cite{DonaldsonKronheimer1990}.
Their success motivates the use of gauge fields in smooth topology,
although it does not settle the
homotopy-sphere case, where the second cohomology vanishes. Twistor methods
also reach beyond the self-dual sector. Witten~\cite{Witten1978} describes
the full classical Yang--Mills equations using holomorphic bundles on the
space of complex null lines, with specified triviality and third-order
formal extension conditions. This is a local
complex flat-spacetime correspondence, not an existence theorem for curved
Riemannian four-manifolds. Together these ideas motivate generalized
coordinate descriptions in which maps depend on both a point and a chiral
direction, and curvature is carried by connection data. For a flat-space
precursor based on generalized conformal maps acting on Yang--Mills fields,
see~\cite{Baker2023} and references therein.

The evolution considered here starts with a smooth Riemannian metric
\(g_0\) on a closed oriented real four-manifold. On a homotopy
four-sphere the conformal condition \(W_g^+=0\) already suffices to
identify the smooth type. To see this, define the chiral Weyl energies
using Hilbert--Schmidt norms on the curvature operators
\(W^\pm_g:\Lambda_g^\pm\to\Lambda_g^\pm\). The signature formula gives
\begin{equation}\label{eq:chiral-weyl-energy}
 E_\pm(g)=\int_M |W^\pm_g|^2_{\mathrm{op,HS}}\,dV_g,
 \qquad E_+(g)-E_-(g)=12\pi^2\tau(M).
\end{equation}
When \(M\) is a homotopy four-sphere, \(\tau(M)=0\), so \(W^+=0\)
forces \(W^-=0\) as well. The metric is then locally conformally flat.
Compactness and simple connectivity make its conformal developing map
to \(S^4\) a diffeomorphism~\cite{Hirzebruch1971Signature,Kuiper1949}.

This suggests evolving the metric by descent for \(E_+\). The
four-index Weyl energy is
\(\mathcal W(g)=\int_M|W_g|_{\mathrm{tensor}}^2\,dV_g=4(E_+(g)+E_-(g))\).
Thus \(E_+=\mathcal W/8+6\pi^2\tau(M)\). On the fixed oriented
four-manifold, their metric gradients differ by a fixed positive factor.
The negative \(L^2\) gradient of \(\mathcal W\) is proportional to the
Bach tensor \(B_g\). Since \(E_+\) is invariant under conformal
rescaling, decreasing it requires an evolution of the conformal
class itself~\cite[\S1 and \S2.3]{BahuaudGuentherIsenbergMazzeo2025Bach}.

We consider the modified Bach flow, the four-dimensional ambient
obstruction flow in the following normalization~\cite[eq.~(1.2)]{BahuaudGuentherIsenbergMazzeo2025Bach}.
\begin{equation}\label{eq:modified-bach}
 \partial_t g=B_g+\frac{1}{12}(\Delta_g s_g)g,\qquad g(0)=g_0 .
\end{equation}
Here \(s_g\) is scalar curvature, \(\Delta_g=\operatorname{tr}_g\nabla^2\),
and \(B_g\) is normalized as in equation~(1.1) of the cited work.
This is fourth order in the metric. The scalar-curvature term addresses
the conformal degeneracy in the short-time existence argument. After
a diffeomorphism gauge adjustment in the proof, the theorem of Bahuaud
and Helliwell gives a smooth short-time solution for every smooth
Riemannian metric \(g_0\) on a closed four-manifold~\cite[Theorem C]{BahuaudHelliwell2011HigherOrder}.
Trace-freeness of \(B_g\) makes the scalar term orthogonal to the
Weyl-energy gradient. Consequently a smooth solution satisfies
\begin{equation}\label{eq:weyl-dissipation}
 \frac{d}{dt}E_+(g(t))=-c_B\int_M |B_g|^2\,dV_g,\qquad c_B>0 ,
\end{equation}
where \(c_B\) accounts for the normalization~\cite[\S2.3]{BahuaudGuentherIsenbergMazzeo2025Bach}.

A single smooth time slice can suffice to recognize the sphere.
Chang, Gursky, and Yang's sphere theorem applies when the Yamabe
constant of the conformal class is positive and
\(\mathcal W(g)<16\pi^2\chi(M)\). On a homotopy four-sphere,
\(\chi(M)=2\) and \(\mathcal W=8E_+\). Thus
\(Y(M,[g])>0\) and \(E_+(g)<4\pi^2\) force the standard smooth
type, since simple connectivity excludes
\(\mathbb{RP}^4\)~\cite[Theorem A]{ChangGurskyYang2003}.

It remains to determine whether the flow can reach this threshold or a
smooth conformally flat limit from broader initial data. The
dissipation identity controls the time integral of the squared Bach
tensor. The continuation criterion of Lopez involves the full Riemann
tensor, and a smooth critical metric satisfies \(B_g=0\).
A global argument must control curvature concentration and show that
the evolution attains a recognition condition~\cite[Theorem 1.2]{Lopez2018AmbientObstruction}.

\bibliographystyle{unsrturl}
\bibliography{references}
\end{document}
